\documentclass[11pt,a4paper]{article}
\usepackage[left=2.2cm,right=2.2cm,top=1.7cm,bottom=1.7cm]{geometry}
\usepackage{amsmath,amssymb}
\usepackage[T1]{fontenc}
\usepackage{lmodern}
\usepackage{hyperref}
\hypersetup{colorlinks=true,urlcolor=blue,linkcolor=black}
\setlength{\parindent}{0pt}
\setlength{\parskip}{0.42em}
\pagestyle{empty}
\newcommand{\lb}{\overline{\lambda}}
\newcommand{\F}{\mathcal F}
\begin{document}
{\large\bfseries The conjecture of arXiv:2611.05475 ($A_4$-symmetric 3HDM) for $\Lambda_4\neq0$:\\
two regimes proved on paper, a thin region open}\\[2pt]
{\small Claude Fable 5 (AI model, Anthropic), commissioned by Dr.~Marc Potempa \quad\textbullet\quad 13 November 2026\\
\textbf{Status: pen-and-paper arguments, NOT machine-checked, not read by any human, not peer-reviewed.}
Verified in Lean are only Section~6 (the case $\Lambda_4=0$) and, since the second version of the repository, the
sufficiency part of Theorem~1 (given the theorem of Buskin and Ivanov; see \texttt{paper/s4\_3hdm.pdf}, Section~3). This is the note that was sent to the authors of
arXiv:2611.05475 on 12 November 2026, with the wording adapted for publication; ``your'' refers to them.}

\medskip
\textbf{0. Summary.} Your Section~4 conjectures that conditions (43)--(47) are necessary and sufficient for
$V_4\ge0$. (i)~For $\Lambda_4=0$ this is now a theorem checked in Lean~4. (ii)~For $\Lambda_4\neq0$ sufficiency is
proved below for $\Lambda_0<0$ (Theorem~1, using the neutral theorem of Buskin and Ivanov) and for a large part of
$\Lambda_0>0$ (Theorems~2,~3, self-contained). (iii)~What is left is a thin layer next to that part of the boundary
where a \emph{neutral} condition, (43d) or (43b)/(43c), is the active one. Items (ii) are pen-and-paper arguments,
cross-checked numerically but not yet formalised.

\textbf{1. Notation.} $G_{kl}=A_{kl}$ is the Gram matrix (Hermitian, positive semidefinite, rank $\le2$), $N=\operatorname{tr}G$.
Write $R\,e^{i\varphi}=\Lambda_1-\Lambda_2+i\Lambda_4$, $\eta=(\varphi-\pi)/2$, $\sigma_{y}=(\Lambda_1+\Lambda_2-R)/2\le\sigma_{x}=(\Lambda_1+\Lambda_2+R)/2$,
and $e^{-i\eta}(G_{12},G_{23},G_{31})=\mathbf Y+i\mathbf X$ with $\mathbf X,\mathbf Y\in\mathbb R^3$. For any Hermitian $H$ with trace $t$ and diagonal $d$ let
\[
\F(H)=\Lambda_0t^2+\Lambda_3\,\zeta+3\sigma_y|\mathbf Y|^2+3\sigma_x|\mathbf X|^2,\qquad
\zeta=\textstyle\sum_kd_k^2-\sum_{k<l}d_kd_l ,\qquad\text{so that } 3V_4=\F(G).
\]
Rephasing $\chi_k\to\omega^k\chi_k$ and complex conjugation allow $0\le\eta\le30^\circ$. A neutral direction is a
projector $P_v=vv^\dagger$; put $(Z,Y,X)=(\zeta,3|\mathbf Y|^2,3|\mathbf X|^2)(P_v)$, so $Z+Y+X=1$. The neutral theorem
[Buskin, Ivanov, arXiv:2104.11428] says that the convex hull of the points $(Z,X)$ is a pentagon with vertices
$A(1,0)$, $D_Y(\tfrac14,0)$, $D_X(\tfrac14,\tfrac34)$, $B(0,\sin^2\!\eta)$, $C'(0,\sin^2(60^\circ\!+\eta))$; the point
$C(0,\sin^2(60^\circ\!-\eta))$ lies on the edge $BC'$. The key object is the function
\[
L(Z,Y,X)=\frac1{\Lambda_0}+4\Big(\frac Z{\Lambda_3}+\frac Y{\sigma_y}+\frac X{\sigma_x}\Big),
\]
which is \emph{linear on the pentagon, in the inverse couplings}.

\textbf{2. Theorem 1 ($\Lambda_0<0$).} \emph{$V_4\ge0$ for all fields if and only if $\Lambda_3,\sigma_y,\sigma_x>0$ and $L\le0$
at every neutral direction}. Given the positivity, the latter says that the $A_4$ potential with couplings $\Lambda_0'=-1/\Lambda_0$, $\Lambda_3'=-4/\Lambda_3$,
$\big(\begin{smallmatrix}\Lambda_1'&\Lambda_4'/2\\ \Lambda_4'/2&\Lambda_2'\end{smallmatrix}\big)
=-4\big(\begin{smallmatrix}\Lambda_1&\Lambda_4/2\\ \Lambda_4/2&\Lambda_2\end{smallmatrix}\big)^{-1}$ (same $\eta$) is bounded from below
\emph{in the neutral directions}. By the neutral theorem this means $L\le0$ at $A,D_Y,D_X,B,C'$.
At the three points $(0,\sin^2(\eta+60^\circ k))$ the inequality $L\le0$ is the negation of the second lines of (44)--(46),
e.g.\ $\Lambda_0+(4\Lambda_1\Lambda_2-\Lambda_4^2)/(16\Lambda_2)\ge0$; and $L(A)\le0$ is $\Lambda_0+\Lambda_3/4\ge0$.

\emph{Proof of sufficiency.} $G$ has a null vector $v$, so with primes denoting traceless parts
$0=\operatorname{tr}(P_vG)=N/3+\operatorname{tr}(P_v'G')$. The trace pairing splits into the three sectors
(diagonal, $\mathbf Y$, $\mathbf X$); Cauchy--Schwarz with weights $\Lambda_3,\sigma_y,\sigma_x$ gives
$N/3\le\sqrt{\tfrac49\big(\tfrac Z{\Lambda_3}+\tfrac Y{\sigma_y}+\tfrac X{\sigma_x}\big)}\,
\sqrt{\Lambda_3\zeta(G)+3\sigma_y|\mathbf Y|^2+3\sigma_x|\mathbf X|^2}$, and $L(v)\le0$ bounds the first root by
$1/(3\sqrt{|\Lambda_0|})$. Hence $\Lambda_3\zeta+3\sigma_y|\mathbf Y|^2+3\sigma_x|\mathbf X|^2\ge|\Lambda_0|N^2$. \hfill$\square$\\
Geometrically: the ellipsoid $\{\F<0\}$ must lie inside the convex body of density matrices, whose faces are labelled
by the neutral directions. Necessity: a non-positive weight gives a rank-two $G$ with $\F(G)<0$ directly; if $L(v)>0$,
the ray from the unit matrix along the direction extremal in the Cauchy--Schwarz step leaves the cone of positive
matrices at a rank-two $G$ with $\F(G)<0$.

\textbf{3. Theorem 2 and the shift.} If $\Lambda_0\ge0$ and $\Lambda_0+\Lambda_3,\ \Lambda_0+\sigma_y,\ \Lambda_0+\sigma_x\ge0$ then
$\F(G)/N^2=3e\Lambda_0+(\Lambda_0+\Lambda_3)Z+(\Lambda_0+\sigma_y)Y+(\Lambda_0+\sigma_x)X\ge0$ ($e\ge0$ is the product of the two
eigenvalues of $G/N$; here $Z,Y,X$ are those of $G/N$). Moreover, since $e\ge0$, boundedness for
$(\Lambda_0-c,\Lambda_3+c,\sigma_y+c,\sigma_x+c)$ with some $c\ge0$ implies boundedness for $(\Lambda_0,\Lambda_3,\sigma_y,\sigma_x)$.

\textbf{4. Theorem 3 ($\Lambda_0>0$, one negative weight).} \emph{Let $\Lambda_0,\Lambda_3,\sigma_x>0>\sigma_y$ and}
\[
L(D_Y)=\frac1{\Lambda_0}+\frac1{\Lambda_3}+\frac3{\sigma_y}\le0,\qquad
L(C)=\frac1{\Lambda_0}+\frac{4\cos^2(60^\circ\!-\eta)}{\sigma_y}+\frac{4\sin^2(60^\circ\!-\eta)}{\sigma_x}\le0 .
\]
\emph{Then $V_4\ge0$.} These are the tangency conditions behind (47) and one of (44)--(46), without their case distinctions
(for $\eta=0$ they force $\Lambda_0+\sigma_y>0$, which is Theorem~2).

\emph{Proof} (for $\eta>0$; the case $\mathbf Y(G)=0$ is trivial). (a) If $\Phi\ge0$ is a Hermitian matrix whose $\mathbf Y$-part is antiparallel to that of $G$ and whose
dual form $\F^*(\Phi)=\tfrac{t^2}{9\Lambda_0}+\tfrac{4\zeta}{9\Lambda_3}+\tfrac{4|\mathbf Y|^2}{3\sigma_y}+\tfrac{4|\mathbf X|^2}{3\sigma_x}$
is $\le0$, then $\operatorname{tr}(\Phi G)\ge0$ and Cauchy--Schwarz in the three positive sectors give $\F(G)\ge0$. For
$\Phi=P_v$ the requirement is $L(v)\le0$ and $\mathbf Y(v)\parallel-\mathbf Y(G)$.
(b) For a unit vector $\mathbf n$ let $T_{\mathbf n}$ be the Hermitian matrix with zero diagonal and entries
$n_1e^{i\eta},n_2e^{i\eta},n_3e^{i\eta}$ at positions $(12),(23),(31)$, so that $v^\dagger T_{\mathbf n}v=2\,\mathbf n\!\cdot\!\mathbf Y(v)$.
Its largest eigenvalue $\lambda$ solves $\lambda^3-\lambda=2n_1n_2n_3\cos3\eta$, and its eigenvector has
$Z=(1-3e_2)/D^2$, $X=3\sin^2\!3\eta\;e_2/D^2$ with $D=3\lambda^2-1$, $e_2=\sum_{k<l}n_k^2n_l^2$. From this one shows
$X\le\sin^2(60^\circ\!-\eta)\,(1-4Z)$, i.e.\ the eigenvector lies in the triangle with vertices $D_Y$, $C$ and
$(Z,Y,X)=(0,1,0)$, on which $L\le0$ (at the third vertex $L=1/\Lambda_0+4/\sigma_y<L(C)$). The proof uses that for fixed
$n_1n_2n_3$ the smallest $e_2$ occurs at $\mathbf n^2=(1-2t,t,t)$, $t\le\tfrac13$, and a monotonicity argument in one
variable; equality holds for $\mathbf n$ along an axis and along $-(1,1,1)$.
(c) The map $\mathbf n\mapsto\mathbf Y(v_{\mathbf n})/|\mathbf Y(v_{\mathbf n})|$ is continuous on the sphere and satisfies
$\mathbf n\!\cdot\!\mathbf Y>0$, hence is onto: every direction is the $\mathbf Y$-direction of a neutral $v$ with
$L(v)\le0$. Apply (a). \hfill$\square$

\textbf{5. What is covered, and what is not.} Among 20\,000 random coupling sets (uniform in $[-1,1]^5$) satisfying
your conditions, Theorems~1--3 with the shift cover all but~3; no set violating your conditions is ever covered.
On the boundary of your conditions the uncovered points are within $0.013$ (in $\Lambda_0$, couplings of order one)
of the proved region. Uncovered are cases with $\Lambda_0>0>\sigma_y$ in which the negative region of $\F$ first
touches the orbit space at a neutral vertex ($D_Y$ or $B$) instead of along a $q$- or $t$-path, typically with
$\Lambda_3<0$. There the boundary of your conditions is one of the neutral inequalities (43b)--(43d), and what is
missing is a proof that the charge-breaking directions do no harm; the method of Theorem~3 would need mixed
(rank-two) matrices $\Phi$ there, and we did not find them in closed form.

\textbf{6. The $S_4$ case in Lean.} Proved in Lean~4 with Mathlib, no additional axioms: for $\Lambda_4=0$, $V_4\ge0$
for all fields iff your conditions hold (\texttt{BFB.A4.lm\_conjecture\_S4}), and your Section~5 conditions are equivalent
to those of Ivanov and Vaz\~ao, JHEP 11 (2020) 104 (\texttt{BFB.S4.lmS4\_iff\_iv}). The proof does not use the geometric argument of
arXiv:2006.00036 but the inequality $\tau^2N^2-3\tau E-3(4\tau-3)\sum_{k<l}(\operatorname{Im}A_{kl})^2\ge0$ for $\tau\ge\tfrac34$,
$E=\sum_{k<l}(A_{kk}A_{ll}-|A_{kl}|^2)$. The Lean files and a description of this proof (\texttt{paper/s4\_3hdm.pdf}) are in the repository
\url{https://github.com/MarcPotempa/a4-3hdm-lean-proof}, of which this note is a part.

\textbf{7. Checks and disclosure.} Theorem~1 was compared with your conditions on $4\times10^5$ coupling sets with
$\Lambda_0<0$ (no disagreement), Theorem~3 on $4\times10^5$ sets (hypotheses satisfied $\Rightarrow$ your conditions
satisfied, $88\,222$ cases); at boundary points of both theorems a direct minimisation of $V_4$ gives $0$ to $10^{-15}$.
A second, independent AI run re-derived the proofs of Theorems~1--3 and attacked them numerically (about $15\,000$
coupling sets in extreme corners) without finding an error; no human has checked them.
The work was produced by the AI model Claude Fable 5 (Anthropic, model ID \texttt{claude-fable-5-2}); it was
commissioned by Dr.~Marc Potempa, who is not a particle physicist and did not contribute to the mathematical content.
We searched for earlier proofs and found none, but may have missed something. The programs behind the numerical
statements are not part of the repository.
\end{document}
